\chapter[Requirements and System Architecture]{Requirements and System Architecture}
\label{chap:requirements-architecture}

This chapter translates the problem introduced in Chapter~\ref{chap:introduction} into concrete
system requirements and presents the architecture used to satisfy them. The
architecture is described before the detailed processing and retrieval
algorithms because those algorithms depend on several system-wide decisions:
documents are owned through workspaces, processing is asynchronous, provenance
is retained from extraction to citation, and text obtained from a document is
treated as untrusted input. The descriptions in this chapter are based on the
implemented application rather than on a hypothetical target design.

The application is organized as a modular monolith. A React single-page
application communicates with an ASP.NET Core JavaScript Object Notation
(JSON) application programming interface (API); the backend coordinates
local document processing, PostgreSQL persistence, model-provider calls, and a
background ingestion worker. This arrangement keeps deployment manageable for
a diploma project while still separating responsibilities through explicit
service and provider interfaces. The separation is important for testing and
for failure isolation, even though the modules execute in one server process.

\section{Functional requirements}
\label{sec:functional-requirements}

Functional requirements state what an authenticated user must be able to
accomplish. They were derived from the central use case: a user creates a
workspace, contributes documents to its reusable knowledge collection, and asks
questions whose answers can be checked against those documents. The word
\emph{workspace} is used in the interface, while \emph{Project} remains the
backend and persistence term. Table~\ref{tab:functional-requirements} summarizes
the resulting requirements and gives an observable acceptance condition for
each one.

{\small
\begin{longtable}{@{}L{1.2cm}L{7.9cm}L{4.6cm}@{}}
\caption[Functional requirements]{Functional requirements and observable acceptance conditions.}
\label{tab:functional-requirements}\\
\toprule
\textbf{ID} & \textbf{Requirement} & \textbf{Acceptance condition} \\
\midrule
\endfirsthead
\toprule
\textbf{ID} & \textbf{Requirement} & \textbf{Acceptance condition} \\
\midrule
\endhead
\bottomrule
\endfoot
\textbf{FR-01} &
The system shall support account registration, sign-in, session checking, and
sign-out. &
Protected API operations reject an unauthenticated request, while a valid
session identifies one application user. \\

\textbf{FR-02} &
The user shall be able to create, list, edit, open, and delete owned workspaces. &
A user sees and changes only projects whose owner identifier matches the
authenticated identity. \\

\textbf{FR-03} &
The user shall be able to upload supported Portable Document Format (PDF) and
Office Open XML word-processing (DOCX) documents to an owned workspace. The
server shall validate the request before accepting the file. &
An accepted file has an allowed extension, matching declared media type,
recognizable file signature, safe original name, and size within the configured
limit. \\

\textbf{FR-04} &
The system shall transform an accepted document into reusable searchable
knowledge while preserving its source order and provenance. &
Successful processing produces raw and normalized source sections, bounded
chunks, embeddings, and aggregate processing metadata, and places the document
in the \emph{Ready} state. \\

\textbf{FR-05} &
The system shall diagnose the technical representation of PDF pages and apply
local optical character recognition (OCR) selectively when a page is classified
as scanned. &
Page diagnostics and OCR outcomes are inspectable, and successfully recognized
text records OCR rather than native extraction as its provenance. \\

\textbf{FR-06} &
The user shall be able to inspect document processing results and request
controlled rebuilds. &
The interface exposes processing, normalization, chunking, embedding, technical
analysis, OCR, and Document Understanding information; rebuild requests are
accepted only in compatible states. \\

\textbf{FR-07} &
The system shall retrieve relevant evidence within one owned workspace by
combining semantic, lexical, and document-metadata signals. &
A search request returns only eligible chunks from \emph{Ready} documents in
the selected project, together with the available rank diagnostics. \\

\textbf{FR-08} &
The system shall answer a question from bounded retrieved evidence and attach
traceable citations. &
An answer either refers to backend-validated source identifiers or states that
the workspace documents do not provide sufficient evidence. \\

\textbf{FR-09} &
The user shall be able to maintain multiple persistent conversations inside a
workspace. &
Conversations can be created, opened, renamed, and deleted; messages retain a
stable order and the first question produces a concise default title. \\

\textbf{FR-10} &
The system shall preserve enough source information to inspect a previously
generated answer. &
Each cited assistant message stores a bounded snapshot containing the source
identifier, document name, chunk and page information, heading, and excerpt. \\
\end{longtable}
}

These requirements deliberately place documents at workspace level rather than
attaching them to a single conversation. Extraction, chunking, and embedding
can therefore be performed once and reused by every conversation in that
workspace. A conversation supplies interaction history, but it does not own or
duplicate the underlying knowledge collection. The same distinction also
supports direct search and diagnostics independently of answer generation.

\section{Non-functional requirements}
\label{sec:non-functional-requirements}

For this application, correctness includes more than returning fluent text. A
response that uses another user's document, loses the location of its evidence,
or silently accepts a stale embedding is incorrect even if its wording appears
plausible. The non-functional requirements in
Table~\ref{tab:nonfunctional-requirements} therefore emphasize isolation,
traceability, predictable resource use, and recoverability. They are
project-specific engineering requirements rather than claims of compliance
with an external certification standard.

{\small
\begin{longtable}{@{}L{1.2cm}L{5.2cm}L{7.3cm}@{}}
\caption[Non-functional requirements]{Non-functional requirements and their architectural responses.}
\label{tab:nonfunctional-requirements}\\
\toprule
\textbf{ID} & \textbf{Quality requirement} & \textbf{Architectural response} \\
\midrule
\endfirsthead
\toprule
\textbf{ID} & \textbf{Quality requirement} & \textbf{Architectural response} \\
\midrule
\endhead
\bottomrule
\endfoot
\textbf{NFR-01} &
Ownership isolation &
Every project-scoped read and mutation includes the authenticated owner in its
database predicate. Retrieval repeats project and owner filtering in each
evidence-producing query. \\

\textbf{NFR-02} &
Evidence traceability &
Source sections retain page or heading information and extraction method;
chunks carry source ranges; answers expose only source identifiers that the
backend can map to retrieved chunks. \\

\textbf{NFR-03} &
Bounded resource use &
Upload size, queue capacity, OCR pages and pixels, chunk size, provider inputs
and outputs, retrieval pools, conversation history, and answer context all have
explicit bounds. \\

\textbf{NFR-04} &
Repeatability and deterministic local stages &
Normalization, PDF classification, OCR routing, chunk construction, query
preparation, and rank fusion use versioned or deterministic rules. Content
hashes identify the source used for derived results. \\

\textbf{NFR-05} &
Failure isolation and retryability &
Main document processing, technical PDF analysis, OCR, and Document
Understanding have distinct states. Optional enrichment failure does not
automatically invalidate useful extracted text, while failed core processing
can be queued again. \\

\textbf{NFR-06} &
Persistence consistency &
New source sections, chunks, and validated embeddings are prepared before they
replace the authoritative retrieval representation in a database transaction.
Uniqueness and confidence constraints reinforce application validation. \\

\textbf{NFR-07} &
Maintainability and testability &
Controllers, orchestration services, deterministic algorithms, storage, and
external-model adapters have narrow responsibilities and are connected through
dependency injection, allowing provider boundaries to be replaced by fakes in
tests. \\

\textbf{NFR-08} &
Responsive interaction &
Lengthy ingestion is delegated to a background worker. The API returns the
document state promptly, and the client can refresh visible progress without
holding the upload request open for the whole pipeline. \\

\textbf{NFR-09} &
Controlled information exposure &
Physical storage paths, generated storage names, vectors, credentials, hidden
prompts, and raw provider responses are excluded from public contracts. Public
errors are bounded and phrased without internal details. \\
\end{longtable}
}

The application does not define a formal throughput or response-time service
level. Its resource limits are safety and usability controls for the current
single-instance implementation, not evidence of large-scale performance.
Likewise, deterministic processing means that identical inputs and
configuration produce the same local transformations; it does not mean that an
external generative model is mathematically deterministic.

\section{High-level system architecture}
\label{sec:high-level-architecture}

Figure~\ref{fig:system-architecture} presents the system as four cooperating
runtime areas. First, the browser executes the React and TypeScript client and
holds only interaction state needed to render the current screen. Second, one
ASP.NET Core host exposes the API, enforces identity and ownership, executes the
application services, serves the published frontend, and runs the document
worker. Third, PostgreSQL stores relational state, generated full-text vectors,
and pgvector embeddings, while an application-controlled directory stores the
original uploaded files under generated names. Fourth, the configured model
provider supplies embeddings, structured document analysis, reranking, and
grounded answer generation. PDF inspection, PDF rendering, OCR, normalization,
and chunking remain local to the server.

\thesisfigureplaceholder{DIA-01}{High-level architecture of the AI Document Assistant, showing the browser, modular backend, persistence, local processing, and configured model-provider boundary.}{fig:system-architecture}

Two end-to-end flows cross these areas. On the document side, an authenticated
upload is validated and stored, then a queued job converts it into an ordered,
normalized, provenance-bearing set of chunks and vectors. On the question
side, the current question is embedded and analyzed, several retrieval channels
produce candidates, rank fusion and optional reranking select evidence, and a
bounded context is submitted for answer generation. The backend then validates
source references before returning or persisting them. These flows share the
same project ownership boundary and the same persisted document representation;
they are not independent subsystems joined only at the interface.

Calling the architecture a modular monolith is significant. The design does
not incur network boundaries between extraction, retrieval, and conversation
modules, and a database transaction can protect related authoritative changes.
At the same time, algorithmic and external dependencies are not embedded
directly in controllers. Interfaces isolate text extraction, storage, OCR,
model access, retrieval channels, and answer construction. The present design
is consequently simpler to run locally than a distributed system, but its
process-local queue and local upload directory also make the current deployment
a single-instance design.

\section{Backend architecture}
\label{sec:backend-architecture}

The backend follows a request--service--persistence structure. API controllers
define authenticated JSON or multipart endpoints, validate request-level
conditions, obtain the user identifier from the authenticated principal, and
translate expected failures into bounded HTTP responses. Application services
then coordinate use cases such as processing a document, searching a project,
answering a question, or persisting a conversation. Entity Framework Core
provides the unit-of-work boundary for domain persistence, while direct,
parameterized PostgreSQL commands are used where vector and full-text search
need database-specific operators. The framework mechanisms used for
controllers, dependency injection, hosted services, and relational mapping
require supporting technical references.
\citationneeded{ASP.NET Core controllers, dependency injection, and hosted services}
\citationneeded{Entity Framework Core and the Npgsql PostgreSQL provider}

Provider adapters form deliberate narrow boundaries around facilities that are
either external or implementation-specific. Examples include the local file
store, PDF and DOCX extractors, the PDF structural analyzer, PDFium page
rendering, Tesseract recognition, tokenization, and the four model-provider
tasks. The orchestration code depends on their application interfaces rather
than constructing these clients inside a use case. This is why tests can supply
deterministic embeddings or no-operation analysis providers without changing
the production workflow.

Document jobs are passed through a bounded in-memory channel to one hosted
worker. Duplicate scheduling of the same document identifier is suppressed
while that identifier is pending or running. For each item, the worker creates
a dependency-injection scope and invokes the processing service. The queue is
explicitly process-local and non-durable: a server restart can discard pending
identifiers, although the persisted \emph{Uploaded} or \emph{Failed} document
can be submitted again. This limitation is preferable to describing the queue
as if it provided durable delivery.

Uploaded bytes are stored below a fixed \texttt{Uploads} root using a newly
generated filename whose extension is restricted to PDF or DOCX. Path
construction verifies that the requested storage name is a plain filename and
that its resolved path remains under that root. Public document contracts
return the safe original filename but never the generated storage name or
physical path. Database state and file deletion are coordinated carefully for
document and project deletion, including state checks that prevent deletion
during active processing.

Finally, a global exception handler logs unexpected server failures and returns
a generic JSON message. Expected provider and processing failures are wrapped
in exceptions with safe public messages. Logging statements generally record
identifiers, counts, durations, status, and exception type while deliberately
omitting document text, vectors, prompts, credentials, and provider payloads.
This design supports diagnostics without making internal exception details part
of the public API.

\section{Frontend architecture}
\label{sec:frontend-architecture}

The client is a React single-page application written in TypeScript and built
with Vite. \citationneeded{React and TypeScript official documentation} A small
route parser maps browser paths to authentication, workspace dashboard,
workspace chat, conversation, and document-management views. Navigation uses
the browser history API and responds to back and forward events. At startup,
the client asks the server for the current session; unauthenticated navigation
is redirected to the login interface, while a valid session enters the
workspace area.

State is kept close to the screen that owns it. The chat workspace manages the
current project, document readiness, conversation list, active messages,
uploads, and composer state. The document screen manages document summaries and
loads detailed understanding, technical-analysis, OCR, extracted-text, and
chunk information when required. Shared UI elements provide confirmation
dialogs, notifications, icons, and loading placeholders. This organization is
more appropriate to the present application size than introducing a separate
global state framework.

A typed API module defines the response shapes and centralizes same-origin
requests. It includes the browser cookie on every call, adds a JSON content
type when appropriate, and converts structured server errors into safe messages
for the interface. These TypeScript declarations improve consistency but do
not create a security boundary: browser state and route checks are considered
untrusted, and every authorization decision is repeated by the backend. In a
published build, ASP.NET Core serves the static assets and falls back to the
single application entry page for client routes.

\section{Persistence and domain model}
\label{sec:persistence-domain-model}

The domain model separates user intent, source material, derived retrieval
artifacts, and conversation records. Figure~\ref{fig:domain-model} shows the
principal relationships without reproducing every database field.

\thesisfigureplaceholder{DIA-02}{Compact domain model for workspace ownership, document-processing artifacts, and persistent conversations.}{fig:domain-model}

An \emph{ApplicationUser} owns zero or more \emph{Project} records. Each
Project is one user-facing workspace and is the common parent of its
\emph{Document} and \emph{Conversation} collections. This placement is central
to reuse: all conversations in the project can search the same ready documents,
whereas content from a different project is outside both the retrieval and
authorization scope. Deleting a user cascades to owned projects, and deleting a
project cascades to its database-resident documents and conversations.

A Document represents the uploaded source and the aggregate state of the main
ingestion pipeline. It records the safe original name, generated storage name,
media type, size, timestamps, current processing status, and summary counts or
errors for extraction, normalization, chunking, and embedding. It does not
contain all derived text in one large field. Instead,
\emph{DocumentTextSection} stores the ordered source representation. A section
contains immutable raw extracted text, an optional normalized derivative, its
page number or structural heading where available, and the extraction method.
Consequently, a reader can distinguish native PDF text, OCR text, and DOCX
text, and a rebuild does not destroy the raw material needed for comparison.

\emph{DocumentChunk} is the bounded retrieval unit derived from one or more
normalized sections. Besides its content and stable order, it carries token and
character counts, page range, heading, and the first and last contributing
section indexes. Its embedding is stored in PostgreSQL through the native
pgvector type, together with model, dimensions, content hash, and generation
time. A generated PostgreSQL full-text vector for the same chunk supports the
lexical retrieval channel. This co-location means that relational ownership
filters, text retrieval, and vector retrieval can operate over the same
document and chunk identifiers.
\citationneeded{PostgreSQL generated full-text search columns and generalized inverted indexes}
\citationneeded{pgvector data types and cosine-distance operators}

Three nullable one-to-one aggregates attach distinct forms of document
analysis. \emph{DocumentUnderstanding} contains the controlled document type,
language and confidence information, detected title and subject, model and
prompt identity, source hash, status, and completion information. Its ordered
\emph{DocumentMetadataEntry} children represent bounded items such as
organizations, people, identifiers, dates, and monetary values. In contrast,
\emph{DocumentTechnicalAnalysis} records the structural classification of a
PDF and owns ordered \emph{DocumentPageTechnicalAnalysis} diagnostics.
\emph{DocumentOcrAnalysis} records OCR configuration, routing identity, and
aggregate outcome, with \emph{DocumentPageOcrResult} children only for routed
candidate pages. These aggregates are separate because their lifecycles and
failure meanings differ; a single overloaded document status would lose that
distinction.

The conversation branch has a different purpose. A Conversation belongs to one
Project and orders its \emph{ConversationMessage} children by sequence.
Messages use only the User and Assistant roles. Sources belong only to
assistant messages. Each \emph{ConversationMessageSource} stores a bounded
snapshot of the validated citation display: source identifier, safe document
name, optional document and chunk identifiers, chunk index, page range,
heading, and excerpt. The optional identifiers are scalar historical values,
not foreign keys to the live Document or DocumentChunk rows. Deleting an
individual source document therefore removes its sections, chunks, full-text
index data, and vector values, but does not erase the provenance shown with an
older answer. Deleting the conversation still cascades normally to its
messages and their snapshots.

Several database constraints support these concepts without attempting to
replace application validation. Source-section and chunk indexes are unique
within a document; message sequences are unique within a conversation; source
indexes are unique within an assistant message; and page diagnostic keys
combine their parent analysis with the page number. Confidence and coverage
ratios have database checks that keep them within their valid interval. These
constraints protect ordering and basic integrity while the services enforce
richer currentness rules, such as matching an embedding hash to the exact
current chunk content.

\section{Authentication and workspace ownership}
\label{sec:authentication-ownership}

Authentication uses ASP.NET Core Identity with globally unique user
identifiers and an application cookie.
\citationneeded{ASP.NET Core Identity and cookie authentication}
The cookie is HTTP-only, uses the \emph{Lax} same-site policy, has a sliding
expiration, and follows the security of the current request; HTTPS redirection
is enabled by the host. API authentication and access-denied events return JSON
status responses rather than redirects to HTML pages. Registration also
requires a unique email, while repeated failed sign-in attempts participate in
the configured lockout policy.

Authorization is primarily an ownership rule. A Project contains the
authenticated user's identifier, and each subordinate resource is reached
through that Project. A typical document query therefore requires all three
conditions: the requested document identifier, the project identifier supplied
by the route, and the owning user identifier reached through the Project.
Conversation operations apply the equivalent
Conversation--Project--Owner predicate. A caller cannot combine a conversation
identifier from one workspace with the route of another. Missing resources,
cross-user identifiers, and project/resource mismatches normally produce the
same not-found response, reducing unnecessary disclosure about their
existence.

Retrieval applies the boundary more than once because it is the path that
selects text for external processing. Before generating a query embedding, the
service verifies that the project belongs to the current user. The vector,
lexical, and metadata database queries also include both project and owner
parameters, together with document readiness conditions. Thus a programming
mistake in candidate combination cannot legitimately turn an unfiltered global
chunk into answer context: the evidence-producing queries themselves are
scoped. Parameterized commands are used for database-specific search rather
than interpolating user query text into SQL.

The frontend's current-workspace display and route guards improve navigation
but are not trusted authorization checks. Likewise, resource identifiers are
not treated as secrets. The server assumes that a caller may discover or alter
an identifier and still requires the ownership predicate. The current model
does not implement workspace sharing or cross-user collaboration; each project
has exactly one owner.

\section{Document-processing architecture}
\label{sec:document-processing-architecture}

The upload request performs only the synchronous work required to accept a
document safely. After validating ownership and the file, it stores the bytes,
creates the Document and its initial independent analysis records, and offers
the document identifier to the background queue. A PDF begins with technical
analysis and OCR marked as not yet analyzed, while those concerns are marked
not applicable for DOCX. Document Understanding begins pending. If the bounded
queue cannot accept the item, the persisted document remains visible and can
be manually queued; the API does not pretend that processing was scheduled
durably.

The worker changes an eligible document from \emph{Uploaded} or \emph{Failed}
to \emph{Processing} and executes one coherent pipeline. Technical PDF analysis
runs before extraction so its page classifications can guide OCR. The relevant
extractor produces ordered raw sections; scanned PDF candidates may have
ineffective native text replaced by locally recognized text. Normalization
creates the conservative retrieval representation. Document Understanding
receives a bounded selection of normalized material, then chunking and
embedding produce the searchable units. Chapter~\ref{chap:document-ingestion} explains each transformation
and its provenance rules in detail.

The final source sections, normalized values, chunks, and validated vectors are
constructed before authoritative replacement. The persistence transaction
removes any previous section and chunk representation, inserts the complete new
one, updates aggregate counts and embedding identity, and marks the document
\emph{Ready}. A failure in extraction, normalization, chunking, or embedding
during ordinary processing clears incomplete authoritative retrieval data and
marks the main document \emph{Failed}. Technical analysis and Document
Understanding, however, persist through their own boundaries. Their failure is
recorded without preventing useful normalized text from being chunked and
retrieved. OCR can similarly be partial when some selected pages fail but
usable native or recognized text remains. This status separation is both a
domain-model decision and a resilience mechanism.

\section{Retrieval and question-answering architecture}
\label{sec:retrieval-answering-architecture}

Question answering begins with the current user question, not with a previous
assistant response. The service verifies project ownership, normalizes and
analyzes the question, and generates one embedding in the same configured space
as the document chunks. PostgreSQL then returns bounded vector candidates and
language-neutral lexical candidates, while ready Document Understanding data
may contribute a document-level metadata signal. Only ready documents in the
owned project are eligible. Weighted Reciprocal Rank Fusion deduplicates and
orders the candidate chunks without directly adding incompatible vector and
lexical score scales. A bounded model-based reranker may reorder the strongest
fused candidates; if that stage times out or returns an invalid result, the
hybrid order remains the fallback.
\citationneeded{Retrieval-Augmented Generation}
\citationneeded{Reciprocal Rank Fusion}
\citationneeded{model-based retrieval reranking}

The final chunks are converted into a token-bounded context. The backend assigns
source labels such as \texttt{S1} and includes the document name, page range,
chunk number, heading, and content under explicit untrusted-data delimiters.
The answer provider is instructed to use only this evidence. After generation,
the backend accepts only citation labels that exist in the supplied context and
removes unknown labels. The model can propose the textual marker
\texttt{[S1]}, but it cannot create the authoritative mapping from that marker
to a document.

Persistent conversations add a second, explicitly non-authoritative context.
A bounded set of recent user and assistant messages may help resolve wording,
but it is separated from document evidence and cannot support a citation. The
new question still drives retrieval. The user message is committed before the
provider call, so a temporary answer failure leaves a visible question that can
be retried rather than a fabricated assistant response. On success, the
assistant text and its validated source snapshots are stored together. The
direct search endpoint exposes vector, lexical, metadata, fused, and reranked
diagnostics for evaluation and debugging, while ordinary answer generation
uses the same retrieval service internally.

\section{Security boundaries}
\label{sec:security-boundaries}

The architecture crosses several trust boundaries, summarized in
Table~\ref{tab:security-boundaries}. Security is consequently implemented as a
set of reinforcing controls rather than as a claim that one prompt or one
authentication attribute solves every risk.

{\small
\begin{longtable}{@{}L{3.1cm}L{6.3cm}L{4.6cm}@{}}
\caption[Security boundaries]{Principal trust boundaries, implemented controls, and residual considerations.}
\label{tab:security-boundaries}\\
\toprule
\textbf{Boundary} & \textbf{Implemented controls} & \textbf{Residual consideration} \\
\midrule
\endfirsthead
\toprule
\textbf{Boundary} & \textbf{Implemented controls} & \textbf{Residual consideration} \\
\midrule
\endhead
\bottomrule
\endfoot
Browser to API &
Cookie authentication, server-side validation, owner-scoped database
predicates, bounded contracts, and JSON error handling. &
Client-side routing is not authorization; production transport, cross-site
request protection, and rate limits require deployment review. \\

Uploaded file to local processing &
Size, filename, extension, media-type, and signature checks; generated storage
names; path-containment validation; format-specific parsers. &
The current validation is not a malware scanner, and document parsers still
receive untrusted bytes. \\

Extracted document text to model tasks &
Explicit untrusted-data delimiters, higher-priority task instructions, no
document-driven tool execution, bounded input, structured schemas where
applicable, allowlists, and local output validation. &
Prompt injection is mitigated but not claimed to be completely solved.
Provider behavior and adversarial inputs remain subjects for evaluation. \\

Server to configured model provider &
Provider calls are made only by server-side adapters; keys remain in server
configuration; response payloads and secrets are omitted from public errors and
diagnostic logs. &
Selected document text leaves the local process for embeddings, understanding,
reranking, and answers. Local OCR does not make the complete pipeline local. \\

Generated answer to authoritative citation &
Source identifiers originate in the bounded backend context; unknown model
citations are discarded; persisted sources are validated snapshots. &
A valid identifier establishes provenance, but answer correctness and whether
the cited passage supports every claim still require evaluation. \\
\end{longtable}
}

Prompt-injection controls are applied at three different model boundaries.
Document Understanding asks only for a strict controlled schema and validates
types, metadata categories, lengths, confidence values, and counts locally.
Reranking assigns opaque candidate identifiers, requests only a structured
ordering, rejects invalid relevance values, and fails open to deterministic
hybrid ranking. Grounded answer generation separates retrieved text from
instructions, denies document-driven commands or secret disclosure, and
validates citations after generation. These controls reduce the authority
available to malicious text, but they do not justify describing arbitrary
document content as safe.
\citationneeded{indirect prompt injection and defense-in-depth guidance for language-model applications}

Resource bounds also serve a security role. The application limits accepted
file size, processing queue capacity, OCR raster size and page count, chunk
length, model input and output sizes, retrieval candidate counts, conversation
history, and final context. Such limits constrain accidental and adversarial
resource consumption, although they are not a substitute for deployment-level
rate limiting and monitoring. Database search commands parameterize question
content, while public API contracts exclude physical paths and vector values.

The boundary analysis also clarifies what is not asserted. The current project
does not claim formal multi-tenant certification, complete prompt-injection
prevention, malware detection, encryption of application data at rest, or a
production-scale durable job system. It does establish concrete ownership,
provenance, validation, and failure-isolation rules within the implemented
scope. The student's engineering contribution lies in making these rules hold
across the complete extraction-to-citation path, rather than merely connecting
a document upload control to a general-purpose language model.
